% sec_circle_dimension_algebra.tex -- circle dimension and phase sampling
\section{Circle Dimension and Phase Sampling}\label{sec:circle-dimension}

\subsection{The algebraic core}

The Cayley chart isolates the circle as the basic visible phase factor.
For finitely generated abelian groups this leads to a natural rank-type
invariant defined on the dual torus part.

\begin{definition}[Circle dimension]\label{def:circle-dimension}
Let $G$ be a finitely generated abelian group and write
\[
\widehat G:=\operatorname{Hom}(G,\TT)
\]
for its Pontryagin dual, endowed with the compact-open topology.  We
define the \emph{circle dimension} of~$G$ by
\[
\cdim(G):=\dim\bigl((\widehat G)^0\bigr),
\]
where $(\widehat G)^0$ denotes the identity component of~$\widehat G$.
\end{definition}

\begin{definition}[Half-circle dimension]\label{def:half-circle-dimension}
Let $M$ be a finitely generated cancellative commutative monoid, and let
$G(M)$ be its Grothendieck group completion.  The \emph{half-circle
dimension} of~$M$ is
\[
\cdim_+(M):=\frac{1}{2}\cdim\bigl(G(M)\bigr).
\]
\end{definition}

\begin{proposition}[Basic computation]\label{prop:circle-dimension-basic}
If $G\simeq \ZZ^r\oplus T$ with $T$ finite, then
\[
\widehat G\simeq \TT^r\times \widehat T
\qquad\text{and}\qquad
\cdim(G)=r=\dim_{\QQ}(G\otimes_{\ZZ}\QQ).
\]
In particular,
\[
\cdim(\ZZ^r)=r,
\qquad
\cdim(T)=0
\qquad\text{for every finite abelian group }T,
\]
and
\[
\cdim_+(\NN^d)=\frac{d}{2}
\qquad(d\ge 1).
\]
\leanverified{paper_cdim_basic_computation_seeds}
\end{proposition}

\begin{proof}
Pontryagin duality sends direct sums to direct products, so
\[
\widehat G\simeq \widehat{\ZZ^r}\times \widehat T \simeq \TT^r\times \widehat T.
\]
Since $\widehat T$ is finite, its identity component is trivial, hence
$(\widehat G)^0\simeq \TT^r$ and $\cdim(G)=r$.

The tensor product $G\otimes_{\ZZ}\QQ$ kills the torsion subgroup~$T$,
so
\[
G\otimes_{\ZZ}\QQ \simeq \QQ^r,
\]
and therefore $\dim_{\QQ}(G\otimes_{\ZZ}\QQ)=r$.  The special cases for
$\ZZ^r$ and finite groups are immediate.  Finally,
$G(\NN^d)\simeq \ZZ^d$, so
\[
\cdim_+(\NN^d)=\frac{1}{2}\cdim(\ZZ^d)=\frac{d}{2}.
\]
\end{proof}

\begin{theorem}[Axiomatic rigidity of circle dimension]%
\label{thm:cdim-axiomatic-rigidity}
Let $f$ be a function from the class of finitely generated abelian groups
to $\NN\cup\{0\}$ satisfying:
\begin{enumerate}[label=(\roman*)]
\item $f(G)=f(H)$ whenever $G\simeq H$;
\item $f(G\oplus H)=f(G)+f(H)$ for all $G,H$;
\item if
      $0\to F\to G\to H\to 0$
      is exact with $F$ finite, then $f(G)=f(H)$;
\item $f(\ZZ)=1$;
\item $f(F)=0$ for every finite abelian group~$F$.
\end{enumerate}
Then $f(G)=\cdim(G)$ for every finitely generated abelian group~$G$.
\leanverified{paper_cdim_axiomatic_rigidity_seeds}
\end{theorem}

\begin{proof}
By the structure theorem, every finitely generated abelian group admits a
decomposition
\[
G\simeq \ZZ^r\oplus T
\]
with $T$ finite.  By (iii) and (v),
\[
f(G)=f(\ZZ^r).
\]
Using (ii) and (iv) gives $f(\ZZ^r)=r$.  Proposition~\ref{prop:circle-dimension-basic}
shows that $\cdim(G)=r$, hence $f(G)=\cdim(G)$.
\end{proof}

\begin{theorem}[Short exact additivity and rank-nullity]%
\label{thm:cdim-short-exact-additivity}
Let
\[
0\longrightarrow A\longrightarrow B\longrightarrow C\longrightarrow 0
\]
be a short exact sequence of finitely generated abelian groups.  Then
\[
\cdim(B)=\cdim(A)+\cdim(C).
\]
Consequently, every homomorphism $\varphi:B\to D$ satisfies
\[
\cdim(B)=\cdim(\ker\varphi)+\cdim(\operatorname{im}\varphi).
\]
\leanverified{paper_cdim_short_exact_additivity_seeds}
\end{theorem}

\begin{proof}
By Proposition~\ref{prop:circle-dimension-basic},
\[
\cdim(G)=\dim_{\QQ}(G\otimes_{\ZZ}\QQ)
\]
for every finitely generated abelian group~$G$.  Since $\QQ$ is a flat
$\ZZ$-module, tensoring the given short exact sequence by~$\QQ$ yields
another short exact sequence
\[
0\longrightarrow A\otimes_{\ZZ}\QQ
\longrightarrow B\otimes_{\ZZ}\QQ
\longrightarrow C\otimes_{\ZZ}\QQ
\longrightarrow 0.
\]
Taking dimensions gives the first identity.  The second follows by
applying the first to
\[
0\longrightarrow \ker\varphi
\longrightarrow B
\longrightarrow \operatorname{im}\varphi
\longrightarrow 0.
\]
\end{proof}

\begin{proposition}[Tensor, Hom, and Ext laws]%
\label{prop:cdim-tensor-hom-ext-laws}
Let $G$ and $H$ be finitely generated abelian groups.  Then:
\begin{enumerate}[label=(\roman*)]
\item
      $\cdim(G\otimes_{\ZZ}H)=\cdim(G)\cdot \cdim(H)$;
\item
      $\cdim(\operatorname{Hom}(G,H))=\cdim(G)\cdot \cdim(H)$;
\item
      $\cdim(\operatorname{Ext}^{1}(G,H))=0$.
\end{enumerate}
\leanverified{paper_cdim_tensor_hom_ext_laws_seeds}
\end{proposition}

\begin{proof}
Write
\[
G\simeq \ZZ^r\oplus T,
\qquad
H\simeq \ZZ^s\oplus E,
\]
with $T$ and $E$ finite.  By
Proposition~\ref{prop:circle-dimension-basic},
$r=\cdim(G)$ and $s=\cdim(H)$.

For (i), the free summand of $G\otimes_{\ZZ}H$ is
\[
\ZZ^r\otimes_{\ZZ}\ZZ^s\simeq \ZZ^{rs},
\]
while all other tensor products are torsion.  Hence
$\cdim(G\otimes_{\ZZ}H)=rs$.

For (ii), one has
\[
\operatorname{Hom}(G,H)\simeq
\operatorname{Hom}(\ZZ^r,H)\oplus \operatorname{Hom}(T,H).
\]
The first summand is isomorphic to $H^r$, whose free rank is $rs$; the
second is finite because $T$ is finite and $H$ is finitely generated.
Thus $\cdim(\operatorname{Hom}(G,H))=rs$.

For (iii), additivity in the first variable gives
\[
\operatorname{Ext}^{1}(G,H)\simeq
\operatorname{Ext}^{1}(\ZZ^r,H)\oplus \operatorname{Ext}^{1}(T,H).
\]
The first term vanishes, while $T$ is a finite direct sum of cyclic
groups $C_n=\ZZ/n\ZZ$, and
\[
\operatorname{Ext}^{1}(C_n,H)\simeq H/nH.
\]
Since $H$ is finitely generated, each quotient $H/nH$ is finite, so
$\operatorname{Ext}^{1}(T,H)$ is finite.  Therefore
$\cdim(\operatorname{Ext}^{1}(G,H))=0$.
\end{proof}

\subsection{Phase sampling and reconstruction}

\begin{definition}[Phase-sampling counts]\label{def:cdim-phase-spectrum}
For a finitely generated abelian group $G$ and an integer $N\ge 1$, let
\[
\mu_N:=\{z\in \TT:\ z^N=1\}.
\]
We define the \emph{phase-sampling count} of~$G$ at level~$N$ by
\[
\Sigma_G(N):=\bigl|\operatorname{Hom}(G,\mu_N)\bigr|.
\]
For $N\ge 2$ we also set
\[
\delta_G(N):=\frac{\log \Sigma_G(N)}{\log N}.
\]
\end{definition}

\begin{proposition}[Sampling counts and finite quotients]%
\label{prop:cdim-phase-spectrum-quotient}
For every finitely generated abelian group~$G$ and every integer
$N\ge 1$,
\[
\Sigma_G(N)=|G/NG|.
\]
\leanverified{paper_cdim_phase_spectrum_quotient_seeds}
\end{proposition}

\begin{proof}
It is enough to treat cyclic groups and then use direct-sum
factorization.  For $\ZZ$, one has
\[
\Sigma_{\ZZ}(N)=|\mu_N|=N=|\ZZ/N\ZZ|.
\]
For $C_m=\ZZ/m\ZZ$, a homomorphism $C_m\to\mu_N$ is determined by the
image of a generator, and this image must be an element of $\mu_N$ of
order dividing~$m$.  Hence
\[
\Sigma_{C_m}(N)=\gcd(m,N).
\]
On the other hand,
\[
|C_m/NC_m|=\gcd(m,N).
\]
Taking products over the cyclic decomposition of~$G$ proves the claim.
\end{proof}

\begin{theorem}[Phase-sampling growth recovers circle dimension]%
\label{thm:cdim-phase-spectrum-limit}
Let $G$ be a finitely generated abelian group.  Then
\[
\lim_{N\to\infty}\delta_G(N)=\cdim(G).
\]
More precisely, if $p$ is prime and $a\to\infty$, then
\[
\delta_G(p^a)=\cdim(G)+O(a^{-1}).
\]
\leanverified{paper_cdim_phase_spectrum_limit_seeds}
\end{theorem}

\begin{proof}
Write $G\simeq \ZZ^r\oplus T$ with $T$ finite.  Then
\[
\Sigma_G(N)=\Sigma_{\ZZ^r}(N)\Sigma_T(N)=N^r\Sigma_T(N).
\]
Moreover
\[
1\le \Sigma_T(N)\le |T|.
\]
Therefore
\[
\delta_G(N)=r+\frac{\log\Sigma_T(N)}{\log N}\longrightarrow r.
\]
By Proposition~\ref{prop:circle-dimension-basic}, $r=\cdim(G)$.
For $N=p^a$ one obtains the quantitative bound
\[
\bigl|\delta_G(p^a)-\cdim(G)\bigr|
\le \frac{\log |T|}{a\log p}.
\]
\end{proof}

\begin{theorem}[Reconstruction from the phase spectrum]%
\label{thm:cdim-phase-spectrum-reconstruction}
Let
\[
G\simeq \ZZ^r\oplus\bigoplus_{p}\bigoplus_{a\ge 1}
\bigl(\ZZ/p^a\ZZ\bigr)^{m_{p,a}}
\]
be a finitely generated abelian group.  Then:
\begin{enumerate}[label=(\roman*)]
\item for every prime~$p$,
      \[
      r=\lim_{a\to\infty}\frac{v_p(\Sigma_G(p^a))}{a};
      \]
\item if one sets
      \[
      E_{G,p}(a):=v_p(\Sigma_G(p^a))-ra
      \qquad(a\ge 1),\qquad E_{G,p}(0):=0,
      \]
      then
      \[
      m_{p,a}
      =\bigl(E_{G,p}(a)-E_{G,p}(a-1)\bigr)
      -\bigl(E_{G,p}(a+1)-E_{G,p}(a)\bigr)
      \qquad(a\ge 1).
      \]
\end{enumerate}
Consequently, the full function $N\mapsto \Sigma_G(N)$ determines~$G$
up to isomorphism.
\leanverified{paper_cdim_phase_spectrum_reconstruction_seeds}
\end{theorem}

\begin{proof}
Fix a prime~$p$ and write the $p$-primary part of the torsion subgroup as
\[
T_p\simeq \bigoplus_{j=1}^{t_p}\ZZ/p^{\lambda_{p,j}}\ZZ,
\qquad
\lambda_{p,1}\ge \cdots\ge \lambda_{p,t_p}\ge 1.
\]
By Proposition~\ref{prop:cdim-phase-spectrum-quotient},
\[
\Sigma_G(p^a)=|G/p^aG|
=p^{ra}\prod_{j=1}^{t_p}p^{\min(\lambda_{p,j},a)}.
\]
Hence
\[
v_p(\Sigma_G(p^a))=ra+\sum_{j=1}^{t_p}\min(\lambda_{p,j},a).
\]
Dividing by~$a$ and letting $a\to\infty$ proves (i).

For (ii), define
\[
\Delta E_{G,p}(a):=E_{G,p}(a)-E_{G,p}(a-1)
\qquad(a\ge 1).
\]
Then
\[
\Delta E_{G,p}(a)
=\#\{j:\lambda_{p,j}\ge a\}.
\]
Therefore
\[
m_{p,a}
=\#\{j:\lambda_{p,j}=a\}
=\Delta E_{G,p}(a)-\Delta E_{G,p}(a+1),
\]
which is the stated formula.  The free rank and all multiplicities
$m_{p,a}$ are thus recovered from the phase spectrum, so the structure
theorem yields the isomorphism type of~$G$.
\end{proof}

\subsection{Operational half-circle dimension}

The Cayley phase coordinate also gives a concrete readout model for the
half-circle dimension of the free commutative monoid~$\NN$.

\begin{theorem}[Operational half-circle dimension of $\NN$]%
\label{thm:cdim-operational-half-circle-dimension-N}
For $b\in\NN$ define
\[
q_b(n):=\left\lfloor 2^b\upsilon(n)\right\rfloor
=\left\lfloor \frac{2^b}{\pi}\arctan(n)\right\rfloor
\qquad(n\in\NN),
\]
and let
\[
N_{\max}(b):=
\max\Bigl\{N\in\NN:\ q_b\ \text{is injective on}\ \{0,1,\dots,N\}\Bigr\}.
\]
Then there exist constants $c,C>0$ such that
\[
c\,2^{b/2}\le N_{\max}(b)\le C\,2^{b/2}
\qquad(b\ge 1).
\]
In particular,
\[
\lim_{b\to\infty}\frac{\log_2 N_{\max}(b)}{b}=\frac{1}{2}
=\cdim_+(\NN).
\]
\end{theorem}

\begin{proof}
Set $u(x):=\pi^{-1}\arctan(x)$ for $x\ge 0$.  Since
\[
u'(x)=\frac{1}{\pi(1+x^2)},
\]
the mean value theorem gives
\[
u(n+1)-u(n)\ge \frac{1}{\pi(1+(n+1)^2)}
\qquad(n\ge 0).
\]
Hence, if
\[
\frac{1}{\pi(1+(N+1)^2)}\ge 2^{-b},
\]
then $q_b(0),\dots,q_b(N)$ are strictly increasing, so
$N_{\max}(b)\ge c\,2^{b/2}$ for some $c>0$.

Conversely, suppose that $q_b$ is injective on $\{0,\dots,N\}$ and set
$n:=\lfloor N/2\rfloor$.  Then the values
$q_b(n),q_b(n+1),\dots,q_b(N)$ are distinct, so
\[
u(N)-u(n)\ge \Bigl(\frac{N}{2}-1\Bigr)2^{-b}.
\]
On the other hand,
\[
u(N)-u(n)
=\frac{1}{\pi}\int_n^N\frac{\dd x}{1+x^2}
\le \frac{1}{\pi}\frac{N-n}{1+n^2}
\le \frac{C_0}{N}
\]
for a universal constant $C_0>0$.  Therefore
\[
\Bigl(\frac{N}{2}-1\Bigr)2^{-b}\le \frac{C_0}{N},
\]
which implies $N\le C\,2^{b/2}$ for some $C>0$.
The limit statement follows immediately, and
$\cdim_+(\NN)=1/2$ was computed in
Proposition~\ref{prop:circle-dimension-basic}.
\end{proof}

\begin{corollary}[The box law for $\NN^d$]%
\label{cor:cdim-operational-half-circle-dimension-Nd}
For $d\ge 1$ let
\[
M_d(b):=\bigl(N_{\max}(b)+1\bigr)^d.
\]
Then
\[
M_d(b)=2^{(\frac{d}{2}+o(1))b}
\qquad(b\to\infty).
\]
Equivalently,
\[
\lim_{b\to\infty}\frac{\log_2 M_d(b)}{b}=\frac{d}{2}
=\cdim_+(\NN^d).
\]
\end{corollary}

\begin{proof}
If $q_b$ is injective on $\{0,\dots,N_{\max}(b)\}$, then the coordinatewise
map
\[
(n_1,\dots,n_d)\longmapsto \bigl(q_b(n_1),\dots,q_b(n_d)\bigr)
\]
is injective on the box $\{0,\dots,N_{\max}(b)\}^d$, which has cardinality
$M_d(b)$.  Conversely, every injective dyadic readout on a box
$\{0,\dots,N\}^d$ restricts to an injective readout on each coordinate
axis, hence $N\le N_{\max}(b)$.  Thus the maximal box cardinality is
exactly $(N_{\max}(b)+1)^d$, and the asymptotic law follows from
Theorem~\ref{thm:cdim-operational-half-circle-dimension-N}.  The final
identity is Proposition~\ref{prop:circle-dimension-basic}.
\end{proof}
